\scanpage{41}
$\pi=\pi^p$ satisfait les conditions
\[
\left\{\begin{aligned}
\pi^{\deleted{p}}(x+y)&=\pi x+\pi y+\sum_{0<i<p}\deleted{\unsure{C_i^p}}\frac{1}{i!(p-i)!}x^iy^{\unsure{i}}\\
\pi(fx)&=f^p\pi x
\end{aligned}\right.
\]
Si on a une application satisfaisant ces conditions, peut on définir les $\pi^{(n)}$ \deleted{?} satisfaisant les conditions \deleted{\unsure{premières}} \inserted{des puiss. div.} $(1.1):(1.5)$?, avec $\pi^{\unsure{(p)}}=\pi$

\deleted{On \unsure{construira|construit} les $\pi_N$ par récurrence, en supposant les $\pi_i$ construits pour $i<N$ et \unsure{de façon à} satisfaire $(1.1):(1.5)$ [les $\unsure{1.4|1.5}$ pour $mn<N$, les $1.5$ pour $m+n<N$].

\uline{Cas} $1$\quad $N\not\equiv0\ (p)$ donc $N=N'+\deleted{\unsure{p}}s$, avec $N'\equiv0\ (p)$, $0<s\leq p-1$.

On aura donc $\ldots$}

On va construire
\[
\pi_n(x)=c_nx^{a_0}\pi(x)^{a_1}\pi^2(x)^{a_2}\cdots\pi^r(x)^{a_r}
\]
où
\[
n=a_0+a_1p+\cdots+a_rp^r
\]
est le développement $p$-adique de $n$, et où \deleted{\unsure{c}}
\[
\deleted{c_n\in\mathbb Z\cap\mathbb Z_p^*}\qquad c_n\in\mathbb Q\cap\mathbb Z_p^*
\]
est un élément de $\mathbb Q$ dont la valuation $p$-adique est nulle.

[ Il serait intéressant de connaître $c_n$ mod $p$, qui est un élément de $\mathbb F_p^*$. Est-ce tjs $1$? ].
\mspage{40}
